% =====================================================================
\chapter{Vectores en el plano y en el espacio}
\label{tema:1}
% =====================================================================

Un vector es, a la vez, una lista de números y una flecha con longitud y dirección. Este tema relaciona las dos caras: las operaciones se definen con coordenadas y se interpretan geométricamente. El producto escalar permite medir longitudes y ángulos; el producto vectorial, áreas y volúmenes. La presentación sigue las sesiones~1 a~7 de la unidad~1 del curso 18.02SC del MIT~\cite{mit1802}, las secciones 2.1--2.4 de OpenStax~\cite{openstax3} y la sección~11.1 de Strang~\cite{strangcalc}.

\section{Vectores y operaciones}

\begin{definicion}[Vectores de $\R^n$]\label{def:vector}
Para $n=2$ o $n=3$, un \emph{vector} de $\R^n$ es una lista ordenada $\vu=(u_1,\dots,u_n)$ de números reales, sus \emph{componentes}. Para $\vu,\vv\in\R^n$ y $\lambda\in\R$ se definen la \emph{suma} y el \emph{producto por un escalar} componente a componente:
\[
\vu+\vv=(u_1+v_1,\dots,u_n+v_n),\qquad \lambda\vu=(\lambda u_1,\dots,\lambda u_n).
\]
El vector nulo es $\vcero=(0,\dots,0)$ y $-\vu=(-1)\vu$. En $\R^3$ se escribe $\vi=(1,0,0)$, $\vj=(0,1,0)$, $\vk=(0,0,1)$, de modo que $(u_1,u_2,u_3)=u_1\vi+u_2\vj+u_3\vk$.
\end{definicion}

Las propiedades de la suma y del producto por escalares (asociativa, conmutativa, distributivas, neutro y opuesto) se heredan de las de los números reales, porque se opera componente a componente. Un \emph{punto} y un \emph{vector} se escriben igual; la diferencia es de interpretación. Si $P$ y $Q$ son puntos, el vector $\overrightarrow{PQ}=Q-P$ es la flecha que va de $P$ a $Q$; dos flechas con la misma longitud, dirección y sentido representan el mismo vector, aunque partan de puntos distintos.

Geométricamente, $\vu+\vv$ es la diagonal del paralelogramo de lados $\vu$ y $\vv$ (o, equivalentemente, la flecha que se obtiene colocando $\vv$ a continuación de $\vu$), y $\lambda\vu$ es una flecha paralela a $\vu$, $\abs{\lambda}$ veces más larga y con el sentido invertido si $\lambda<0$.

\begin{ejemplo}\label{ej:t1-punto-medio}
Sean $P=(1,3,-1)$ y $Q=(0,1,1)$. Entonces $\overrightarrow{PQ}=(-1,-2,2)$. El punto medio del segmento $PQ$ es $M=P+\tfrac12\overrightarrow{PQ}=\tfrac12(P+Q)=(\tfrac12,2,0)$: el punto medio es la media de los extremos.
\end{ejemplo}

\section{Producto escalar, norma y ángulo}

\begin{definicion}[Producto escalar y norma]\label{def:escalar}
El \emph{producto escalar} de $\vu,\vv\in\R^n$ es el número
\[
\vu\cdot\vv=u_1v_1+\dots+u_nv_n=\sum_{i=1}^{n}u_iv_i .
\]
La \emph{norma} (o longitud) de $\vu$ es $\norm{\vu}=\sqrt{\vu\cdot\vu}=\sqrt{u_1^2+\dots+u_n^2}$, y la \emph{distancia} entre $\vu$ y $\vv$ es $\norm{\vu-\vv}$. Un vector de norma~1 es \emph{unitario}; si $\vu\neq\vcero$, el vector $\vu/\norm{\vu}$ es el unitario en la dirección de $\vu$.
\end{definicion}

En $\R^2$ y $\R^3$ la norma es la longitud de la flecha, por el teorema de Pitágoras (aplicado dos veces en $\R^3$).

\begin{proposicion}[Propiedades del producto escalar]\label{prop:escalar}
Para todos $\vu,\vv,\vw\in\R^n$ y $\lambda\in\R$:
\begin{enumerate}
  \item simetría: $\vu\cdot\vv=\vv\cdot\vu$;
  \item linealidad: $(\vu+\vw)\cdot\vv=\vu\cdot\vv+\vw\cdot\vv$ y $(\lambda\vu)\cdot\vv=\lambda(\vu\cdot\vv)$;
  \item positividad: $\vu\cdot\vu\ge 0$, con igualdad si y solo si $\vu=\vcero$;
  \item $\norm{\lambda\vu}=\abs{\lambda}\,\norm{\vu}$.
\end{enumerate}
\end{proposicion}

\begin{proof}
(i) y (ii) se siguen de las propiedades de los números reales aplicadas a cada sumando: por ejemplo, $\sum(u_i+w_i)v_i=\sum u_iv_i+\sum w_iv_i$. (iii) $\vu\cdot\vu=\sum u_i^2$ es una suma de cuadrados, no negativa, y vale~0 solo si todos los $u_i$ son nulos. (iv) $\norm{\lambda\vu}=\sqrt{\sum\lambda^2u_i^2}=\sqrt{\lambda^2}\sqrt{\sum u_i^2}=\abs{\lambda}\norm{\vu}$.
\end{proof}

Desarrollando con (i) y (ii) se obtiene la identidad, que se usará muchas veces,
\begin{equation}\label{eq:t1-desarrollo}
\norm{\vu\pm\vv}^2=(\vu\pm\vv)\cdot(\vu\pm\vv)=\norm{\vu}^2\pm2\,\vu\cdot\vv+\norm{\vv}^2 .
\end{equation}

\begin{teorema}[Desigualdad de Cauchy--Schwarz]\label{teo:cs}
Para todos $\vu,\vv\in\R^n$, $\abs{\vu\cdot\vv}\le\norm{\vu}\,\norm{\vv}$. La igualdad se da si y solo si uno de los vectores es múltiplo del otro.
\end{teorema}

\begin{proof}
Si $\vv=\vcero$ ambos miembros son nulos y $\vv=0\cdot\vu$. Si $\vv\neq\vcero$, sea $t=(\vu\cdot\vv)/\norm{\vv}^2$. Por la positividad y por~\eqref{eq:t1-desarrollo} con $t\vv$ en lugar de $\vv$,
\[
0\le\norm{\vu-t\vv}^2=\norm{\vu}^2-2t\,\vu\cdot\vv+t^2\norm{\vv}^2=\norm{\vu}^2-\frac{(\vu\cdot\vv)^2}{\norm{\vv}^2}.
\]
Multiplicando por $\norm{\vv}^2>0$, $(\vu\cdot\vv)^2\le\norm{\vu}^2\norm{\vv}^2$, y tomando raíces se obtiene la desigualdad. La igualdad equivale a $\norm{\vu-t\vv}=0$, es decir, a $\vu=t\vv$. Recíprocamente, si $\vu=s\vv$, ambos miembros valen $\abs{s}\norm{\vv}^2$.
\end{proof}

\begin{corolario}[Desigualdad triangular]\label{cor:triangular}
Para todos $\vu,\vv\in\R^n$, $\norm{\vu+\vv}\le\norm{\vu}+\norm{\vv}$.
\end{corolario}

\begin{proof}
Por~\eqref{eq:t1-desarrollo} y Cauchy--Schwarz, $\norm{\vu+\vv}^2=\norm{\vu}^2+2\,\vu\cdot\vv+\norm{\vv}^2\le\norm{\vu}^2+2\norm{\vu}\norm{\vv}+\norm{\vv}^2=(\norm{\vu}+\norm{\vv})^2$. Ambos miembros son no negativos, luego se pueden tomar raíces.
\end{proof}

Por Cauchy--Schwarz, si $\vu,\vv\neq\vcero$ el cociente $\vu\cdot\vv/(\norm{\vu}\norm{\vv})$ está en $[-1,1]$, lo que justifica la definición siguiente.

\begin{definicion}[Ángulo y ortogonalidad]\label{def:angulo}
El \emph{ángulo} entre dos vectores no nulos $\vu$ y $\vv$ es el único $\theta\in[0,\pi]$ tal que
\[
\cos\theta=\frac{\vu\cdot\vv}{\norm{\vu}\,\norm{\vv}} .
\]
Dos vectores son \emph{ortogonales} (o perpendiculares), $\vu\perp\vv$, si $\vu\cdot\vv=0$.
\end{definicion}

\begin{proposicion}[El ángulo algebraico es el geométrico]\label{prop:coseno}
Si $\vu,\vv\neq\vcero$ forman, como flechas con origen común, un ángulo geométrico $\theta$, entonces $\vu\cdot\vv=\norm{\vu}\norm{\vv}\cos\theta$.
\end{proposicion}

\begin{proof}
El triángulo de lados $\vu$, $\vv$ y $\vu-\vv$ tiene lados de longitudes $\norm{\vu}$, $\norm{\vv}$ y $\norm{\vu-\vv}$, y el ángulo comprendido entre los dos primeros es $\theta$. El teorema del coseno de la trigonometría plana (módulo~0.4) da $\norm{\vu-\vv}^2=\norm{\vu}^2+\norm{\vv}^2-2\norm{\vu}\norm{\vv}\cos\theta$. Comparando con~\eqref{eq:t1-desarrollo} con el signo menos, $\vu\cdot\vv=\norm{\vu}\norm{\vv}\cos\theta$. Si $\vu$ y $\vv$ son paralelos, el triángulo degenera y la igualdad se comprueba directamente ($\theta=0$ o $\theta=\pi$).
\end{proof}

Así, el signo de $\vu\cdot\vv$ informa del ángulo: positivo si es agudo, nulo si es recto y negativo si es obtuso. Si $\vu\perp\vv$, la identidad~\eqref{eq:t1-desarrollo} se reduce al \emph{teorema de Pitágoras} $\norm{\vu+\vv}^2=\norm{\vu}^2+\norm{\vv}^2$.

\begin{ejemplo}\label{ej:t1-angulo}
Sean $\vu=(1,2,2)$ y $\vv=(2,0,1)$. Entonces $\vu\cdot\vv=2+0+2=4$, $\norm{\vu}=\sqrt{1+4+4}=3$ y $\norm{\vv}=\sqrt5$, de modo que $\cos\theta=4/(3\sqrt5)\approx0{,}5963$ y $\theta\approx53{,}40^\circ$. Cauchy--Schwarz se cumple: $4\le3\sqrt5\approx6{,}708$.
\end{ejemplo}

\section{Proyecciones}

\begin{definicion}[Proyección ortogonal]\label{def:proyeccion}
Sean $\vu,\vv\in\R^n$ con $\vv\neq\vcero$. La \emph{componente escalar} de $\vu$ en la dirección de $\vv$ es $\operatorname{comp}_{\vv}\vu=\vu\cdot\vv/\norm{\vv}$, y la \emph{proyección ortogonal} de $\vu$ sobre $\vv$ es el vector
\[
\proj_{\vv}\vu=\frac{\vu\cdot\vv}{\vv\cdot\vv}\,\vv .
\]
\end{definicion}

\begin{proposicion}[Descomposición ortogonal]\label{prop:descomposicion}
Sea $\vv\neq\vcero$. Todo $\vu$ se escribe de forma única como $\vu=\vu_\parallel+\vu_\perp$ con $\vu_\parallel$ múltiplo de $\vv$ y $\vu_\perp\perp\vv$; además, $\vu_\parallel=\proj_{\vv}\vu$. Entre todos los múltiplos $t\vv$, el más cercano a $\vu$ es $\proj_{\vv}\vu$.
\end{proposicion}

\begin{proof}
Si $\vu=t\vv+\vu_\perp$ con $\vu_\perp\cdot\vv=0$, multiplicando escalarmente por $\vv$ se obtiene $\vu\cdot\vv=t\,\vv\cdot\vv$, luego $t=\vu\cdot\vv/\vv\cdot\vv$ está determinado y $\vu_\parallel=\proj_{\vv}\vu$. Recíprocamente, con ese $t$, el vector $\vu_\perp=\vu-t\vv$ cumple $\vu_\perp\cdot\vv=\vu\cdot\vv-t\,\vv\cdot\vv=0$. Para la última afirmación, para cualquier $s\in\R$ el vector $\vu-s\vv=\vu_\perp+(t-s)\vv$ es suma de dos vectores ortogonales, y por Pitágoras $\norm{\vu-s\vv}^2=\norm{\vu_\perp}^2+(t-s)^2\norm{\vv}^2$, que es mínimo exactamente cuando $s=t$.
\end{proof}

\begin{figure}[ht]
\centering
\begin{tikzpicture}[>={Stealth[length=2mm]},scale=1.1,font=\small]
  \draw[->,thick] (0,0) -- (5,0) node[below right] {$\vv$};
  \draw[->,thick] (0,0) -- (3,2.2) node[above] {$\vu$};
  \draw[->,very thick,gray] (0,0) -- (3,0) node[below,black] {$\proj_{\vv}\vu$};
  \draw[->,dashed] (3,0) -- (3,2.2) node[midway,right] {$\vu_\perp$};
  \draw (2.8,0) -- (2.8,0.2) -- (3,0.2);
  \draw (0.8,0) arc (0:36:0.8); \node at (1.05,0.3) {$\theta$};
\end{tikzpicture}
\caption{Proyección ortogonal de $\vu$ sobre $\vv$. Su longitud con signo es $\norm{\vu}\cos\theta=\operatorname{comp}_{\vv}\vu$.}
\label{fig:t1-proyeccion}
\end{figure}

\begin{ejemplo}\label{ej:t1-proyeccion}
Con $\vu=(3,4)$ y $\vv=(1,0)$, $\proj_{\vv}\vu=(3,0)$ y $\vu_\perp=(0,4)$: proyectar sobre un eje es quedarse con la coordenada correspondiente. Con $\vv=(1,1)$, en cambio, $\proj_{\vv}\vu=\tfrac{7}{2}(1,1)=(\tfrac72,\tfrac72)$ y $\vu_\perp=(-\tfrac12,\tfrac12)$, que efectivamente cumple $\vu_\perp\cdot\vv=0$.
\end{ejemplo}

\section{Producto vectorial y producto mixto}

\begin{definicion}[Producto vectorial]\label{def:vectorial}
El \emph{producto vectorial} de $\vu,\vv\in\R^3$ es el vector
\[
\vu\times\vv=(u_2v_3-u_3v_2,\ u_3v_1-u_1v_3,\ u_1v_2-u_2v_1)
=\begin{vmatrix}\vi&\vj&\vk\\u_1&u_2&u_3\\v_1&v_2&v_3\end{vmatrix},
\]
donde el determinante es simbólico y se desarrolla por la primera fila (\cref{def:det3}).
\end{definicion}

\begin{proposicion}[Propiedades del producto vectorial]\label{prop:vectorial}
Para todos $\vu,\vv,\vw\in\R^3$ y $\lambda\in\R$:
\begin{enumerate}
  \item anticonmutatividad: $\vv\times\vu=-\vu\times\vv$; en particular, $\vu\times\vu=\vcero$;
  \item linealidad: $(\vu+\vw)\times\vv=\vu\times\vv+\vw\times\vv$ y $(\lambda\vu)\times\vv=\lambda(\vu\times\vv)$;
  \item ortogonalidad: $(\vu\times\vv)\cdot\vu=(\vu\times\vv)\cdot\vv=0$;
  \item $\norm{\vu\times\vv}=\norm{\vu}\norm{\vv}\sin\theta$, donde $\theta$ es el ángulo entre $\vu$ y $\vv$; por tanto, $\norm{\vu\times\vv}$ es el área del paralelogramo de lados $\vu$ y $\vv$;
  \item $\vi\times\vj=\vk$, $\vj\times\vk=\vi$, $\vk\times\vi=\vj$.
\end{enumerate}
\end{proposicion}

\begin{proof}
(i), (ii) y (v) se comprueban directamente con la fórmula. (iii) $(\vu\times\vv)\cdot\vu=(u_2v_3-u_3v_2)u_1+(u_3v_1-u_1v_3)u_2+(u_1v_2-u_2v_1)u_3$; los seis productos se cancelan dos a dos. Lo mismo ocurre con $\vv$. (iv) La identidad de Lagrange $\norm{\vu\times\vv}^2=\norm{\vu}^2\norm{\vv}^2-(\vu\cdot\vv)^2$ (\cref{eje:1-lagrange-id}) y la \cref{prop:coseno} dan $\norm{\vu\times\vv}^2=\norm{\vu}^2\norm{\vv}^2(1-\cos^2\theta)=\norm{\vu}^2\norm{\vv}^2\sin^2\theta$, y $\sin\theta\ge0$ en $[0,\pi]$. El paralelogramo tiene base $\norm{\vu}$ y altura $\norm{\vv}\sin\theta$.
\end{proof}

La dirección de $\vu\times\vv$ es, pues, perpendicular al plano que contienen $\vu$ y $\vv$, y su sentido viene dado por la \emph{regla de la mano derecha}: $(\vu,\vv,\vu\times\vv)$ está orientado como $(\vi,\vj,\vk)$. El producto vectorial no es asociativo: $(\vi\times\vi)\times\vj=\vcero$, mientras que $\vi\times(\vi\times\vj)=\vi\times\vk=-\vj$.

\begin{definicion}[Producto mixto]\label{def:mixto}
El \emph{producto mixto} de $\vu,\vv,\vw\in\R^3$ es el número $\vu\cdot(\vv\times\vw)$.
\end{definicion}

\begin{proposicion}[Volumen]\label{prop:volumen}
$\abs{\vu\cdot(\vv\times\vw)}$ es el volumen del paralelepípedo de aristas $\vu$, $\vv$ y $\vw$ con un vértice común. En particular, el producto mixto es nulo si y solo si los tres vectores están en un mismo plano por el origen.
\end{proposicion}

\begin{proof}
La base es el paralelogramo de lados $\vv$ y $\vw$, de área $\norm{\vv\times\vw}$. La altura es la longitud de la proyección de $\vu$ sobre la dirección normal $\vv\times\vw$, es decir, $\abs{\vu\cdot(\vv\times\vw)}/\norm{\vv\times\vw}$ (\cref{def:proyeccion}). Volumen $=$ base $\times$ altura $=\abs{\vu\cdot(\vv\times\vw)}$. (Si $\vv\times\vw=\vcero$, los tres vectores son coplanarios y ambos miembros son nulos.)
\end{proof}

\begin{ejemplo}\label{ej:t1-volumen}
Para $\vu=(1,0,0)$, $\vv=(1,2,0)$, $\vw=(0,1,3)$: $\vv\times\vw=(2\cdot3-0\cdot1,\ 0\cdot0-1\cdot3,\ 1\cdot1-2\cdot0)=(6,-3,1)$ y $\vu\cdot(\vv\times\vw)=6$. El paralelepípedo tiene volumen~6 y el tetraedro con esas tres aristas, la sexta parte: volumen~1.
\end{ejemplo}

\begin{observacion}[Errores frecuentes]
\begin{enumerate}[label=(\alph*)]
\item El producto escalar es un \emph{número}; el vectorial, un \emph{vector}. Expresiones como $(\vu\cdot\vv)\cdot\vw$ no tienen sentido.
\item En general $\norm{\vu+\vv}\neq\norm{\vu}+\norm{\vv}$; la igualdad exige que $\vu$ y $\vv$ tengan la misma dirección y sentido.
\item De $\vu\cdot\vv=\vu\cdot\vw$ con $\vu\neq\vcero$ no se deduce $\vv=\vw$, sino solo $\vu\perp(\vv-\vw)$: no se puede «dividir por un vector».
\item El orden importa en el producto vectorial: $\vu\times\vv=-\vv\times\vu$.
\item Al aplicar $\arccos$ se obtiene un ángulo en $[0,\pi]$; un coseno negativo indica un ángulo obtuso, no un error.
\end{enumerate}
\end{observacion}

\begin{observacion}[Conexión con el curso]
Los \emph{embeddings} de los tokens son vectores de $\R^d$ con $d$ del orden de miles, y su parecido se mide con la \emph{similitud coseno} $\vu\cdot\vv/(\norm{\vu}\norm{\vv})$, que es exactamente el coseno de la \cref{def:angulo} (módulo~2, temas~1 y~2). El mecanismo de atención compara una consulta $\mathbf{q}$ con cada clave $\mathbf{k}$ mediante el producto escalar $\mathbf{q}\cdot\mathbf{k}/\sqrt d$ (módulo~13, tema~2), y Cauchy--Schwarz acota esa puntuación por $\norm{\mathbf{q}}\norm{\mathbf{k}}/\sqrt d$. Las definiciones de este tema valen sin cambios en $\R^d$; solo el producto vectorial es exclusivo de $\R^3$.
\end{observacion}

\begin{lecturas}
\item MIT 18.02SC~\cite{mit1802}, unidad~1, parte~A, sesiones~1--4 (vectores, producto escalar, longitudes y ángulos, componentes) y~7 (producto vectorial).
\item OpenStax \emph{Calculus}, vol.~3~\cite{openstax3}, §2.1 «Vectors in the Plane», §2.2 «Vectors in Three Dimensions», §2.3 «The Dot Product» y §2.4 «The Cross Product».
\item Strang, \emph{Calculus}~\cite{strangcalc}, §11.1 «Vectors and Dot Products» y §11.3 «Cross Products and Determinants».
\end{lecturas}

\section*{Ejercicios del tema 1}
\addcontentsline{toc}{section}{Ejercicios del tema 1}

\begin{ejercicio}{1}{eje:1-vf}
Razone si son verdaderas o falsas, para vectores de $\R^3$: (a) $\vu\cdot\vv$ es un vector; (b) si $\vu\cdot\vv=0$ y $\vu\neq\vcero$, entonces $\vv=\vcero$; (c) $\norm{\vu+\vv}=\norm{\vu}+\norm{\vv}$ siempre; (d) $\vu\times\vv=\vv\times\vu$; (e) $\vu\times\vv$ es ortogonal a $\vu$ y a $\vv$.
\end{ejercicio}

\begin{ejercicio}{1}{eje:1-cosenos}
Sea $\va\neq\vcero$ un vector de $\R^3$ que forma ángulos $\alpha,\beta,\gamma$ con $\vi$, $\vj$ y $\vk$ (sus \emph{ángulos directores}). (a) Demuestre que $\va/\norm{\va}=(\cos\alpha,\cos\beta,\cos\gamma)$ y deduzca que $\cos^2\alpha+\cos^2\beta+\cos^2\gamma=1$. (b) Calcule los cosenos y los ángulos directores de $-\vi+2\vj+2\vk$.
\fuente{adaptado de MIT 18.02SC~\cite{mit1802}, problemas suplementarios, 1A-8.}
\end{ejercicio}

\begin{ejercicio}{2}{eje:1-operaciones}
Sean $\vu=(2,-1,3)$ y $\vv=(1,4,-2)$. Calcule $\vu+\vv$, $3\vu-2\vv$, $\vu\cdot\vv$, $\norm{\vu}$, $\norm{\vv}$ y la distancia entre $\vu$ y $\vv$. ¿Es agudo, recto u obtuso el ángulo que forman?
\end{ejercicio}

\begin{ejercicio}{2}{eje:1-unitarios}
(a) Halle la norma y el vector unitario en la dirección de $\vi+\vj+\vk$, de $2\vi-\vj+2\vk$ y de $3\vi-6\vj-2\vk$. (b) Un vector de norma~6 tiene la dirección de $(\vi+2\vj-2\vk)/3$ y su origen en $(-2,0,1)$. ¿Dónde está su extremo?
\fuente{adaptado de MIT 18.02SC~\cite{mit1802}, problemas suplementarios, 1A-1 y 1A-3.}
\end{ejercicio}

\begin{ejercicio}{2}{eje:1-angulos}
Calcule el ángulo que forman: (a) $\vi-\vk$ y $4\vi+4\vj-2\vk$; (b) $\vi+\vj+2\vk$ y $2\vi-\vj+\vk$.
\fuente{adaptado de MIT 18.02SC~\cite{mit1802}, problemas suplementarios, 1B-1.}
\end{ejercicio}

\begin{ejercicio}{2}{eje:1-ortogonal-c}
¿Para qué valores de $c$ son los vectores $c\vi+2\vj-\vk$ e $\vi-\vj+2\vk$ (a) ortogonales; (b) tales que formen un ángulo agudo?
\fuente{adaptado de MIT 18.02SC~\cite{mit1802}, problemas suplementarios, 1B-2.}
\end{ejercicio}

\begin{ejercicio}{2}{eje:1-proyeccion}
Sean $\mathbf F=2\vi-2\vj+\vk$ y $\va=3\vi+2\vj-6\vk$. Calcule la componente escalar de $\mathbf F$ en la dirección de $\va$, la proyección $\proj_{\va}\mathbf F$ y la descomposición $\mathbf F=\mathbf F_\parallel+\mathbf F_\perp$ de la \cref{prop:descomposicion}. Compruebe que $\mathbf F_\perp\perp\va$.
\fuente{adaptado de MIT 18.02SC~\cite{mit1802}, problemas suplementarios, 1B-5.}
\end{ejercicio}

\begin{ejercicio}{2}{eje:1-cubo}
Sea un cubo de arista~1 con un vértice $P$ en el origen y aristas sobre los ejes. Calcule el ángulo que forma la diagonal principal que parte de $P$ con (a) una diagonal de cara que parte de $P$; (b) una arista que parte de $P$.
\fuente{adaptado de MIT 18.02SC~\cite{mit1802}, problemas suplementarios, 1B-3.}
\end{ejercicio}

\begin{ejercicio}{2}{eje:1-vectorial}
(a) Calcule $\va\times\vb$ para $\va=\vi-2\vj+\vk$ y $\vb=2\vi-\vj-\vk$, y compruebe que el resultado es ortogonal a ambos. (b) Halle el área del triángulo de vértices $P=(2,0,1)$, $Q=(3,1,0)$ y $R=(-1,1,-1)$.
\fuente{adaptado de MIT 18.02SC~\cite{mit1802}, problemas suplementarios, 1D-1 y 1D-2.}
\end{ejercicio}

\begin{ejercicio}{2}{eje:1-tetraedro}
Halle el volumen del tetraedro de vértices $P=(1,0,1)$, $Q=(-1,1,2)$, $R=(0,0,2)$ y $S=(3,1,-1)$. (Indicación: es la sexta parte del paralelepípedo de aristas $\overrightarrow{PQ}$, $\overrightarrow{PR}$, $\overrightarrow{PS}$.)
\fuente{adaptado de MIT 18.02SC~\cite{mit1802}, problemas suplementarios, 1D-7.}
\end{ejercicio}

\begin{ejercicio}{3}{eje:1-polarizacion}
(a) Demuestre, usando solo las propiedades del producto escalar y no las coordenadas, la \emph{identidad de polarización} $\vu\cdot\vv=\tfrac14\norm{\vu+\vv}^2-\tfrac14\norm{\vu-\vv}^2$. (b) Sean $\vu,\vv$ no nulos y no de sentidos opuestos, y $\hat\vu=\vu/\norm{\vu}$, $\hat\vv=\vv/\norm{\vv}$. Demuestre que $\vw=\hat\vu+\hat\vv$ forma el mismo ángulo con $\vu$ que con $\vv$, de modo que $\vw/\norm{\vw}$ es el vector unitario bisector.
\fuente{adaptado de MIT 18.02SC~\cite{mit1802}, \emph{problem set}~1, parte~II, problema~2.}
\end{ejercicio}

\begin{ejercicio}{3}{eje:1-cs-cuadratica}
Dé una segunda demostración de Cauchy--Schwarz: (a) justifique que $p(t)=\norm{\vu+t\vv}^2\ge0$ para todo $t\in\R$; (b) escriba $p(t)$ como un polinomio de segundo grado en $t$ (con $\vv\neq\vcero$); (c) deduzca la desigualdad del discriminante de $p$. (d) Explique, en términos de $p$, cuándo se da la igualdad.
\fuente{adaptado de MIT 18.02SC~\cite{mit1802}, problemas suplementarios, 1B-15.}
\end{ejercicio}

\begin{ejercicio}{3}{eje:1-rombo}
Usando solo vectores, demuestre que las diagonales de un paralelogramo son perpendiculares si y solo si el paralelogramo es un rombo (sus cuatro lados son iguales).
\fuente{adaptado de MIT 18.02SC~\cite{mit1802}, problemas suplementarios, 1B-11.}
\end{ejercicio}

\begin{ejercicio}{3}{eje:1-semicirculo}
Usando solo vectores, demuestre que todo ángulo inscrito en una semicircunferencia es recto: si $A$ y $B$ son extremos de un diámetro de una circunferencia de centro $O$ y $C$ es otro punto de ella, entonces $\overrightarrow{CA}\perp\overrightarrow{CB}$.
\fuente{adaptado de MIT 18.02SC~\cite{mit1802}, problemas suplementarios, 1B-12.}
\end{ejercicio}

\begin{ejercicio}{3}{eje:1-lagrange-id}
(a) Demuestre la \emph{identidad de Lagrange} $\norm{\vu\times\vv}^2=\norm{\vu}^2\norm{\vv}^2-(\vu\cdot\vv)^2$ para $\vu,\vv\in\R^3$ desarrollando ambos miembros en coordenadas. (b) Deduzca que $\vu\times\vv=\vcero$ si y solo si $\vu$ y $\vv$ son paralelos (uno es múltiplo del otro).
\end{ejercicio}
